\documentclass[11pt]{article}

\usepackage[letterpaper,left=1in,right=1in,top=0.75in,bottom=0.75in]{geometry}
\usepackage{microtype}
\usepackage{amsmath}
\usepackage{amssymb}
\usepackage{booktabs}
\usepackage{array}
\usepackage{tabularx}
\usepackage[table,dvipsnames]{xcolor}
\usepackage{enumitem}
\usepackage[most]{tcolorbox}
\usepackage{float}
\usepackage[hidelinks]{hyperref}

\emergencystretch=3em

\definecolor{infoblue}{HTML}{2C3E50}
\definecolor{rule}{HTML}{BDC3C7}
\definecolor{decisiongreen}{HTML}{1B7F5A}
\definecolor{boundaryred}{HTML}{A33A2B}

\newcommand{\code}[1]{\texttt{#1}}
\newcommand{\rfckw}[1]{\textsc{#1}}
\newcommand{\painpaper}{\href{https://stackarmor.github.io/rfc-fedramp-vdr/vdr-pain-cvss.pdf}{the PAIN method}}
\newcolumntype{P}[1]{>{\raggedright\arraybackslash}p{#1}}
\newcolumntype{Y}{>{\raggedright\arraybackslash}X}

\newtcolorbox{plain}{
  enhanced, breakable,
  colback=Sepia!7, colframe=Sepia!55, boxrule=0.4pt, arc=2pt,
  left=7pt, right=7pt, top=4pt, bottom=4pt,
  fontupper=\small,
  coltitle=Sepia!72!black, fonttitle=\bfseries\sffamily\footnotesize,
  title={IN PLAIN TERMS}
}

\newtcolorbox{modelbox}{
  enhanced, breakable,
  colback=decisiongreen!5, colframe=decisiongreen!65, boxrule=0.5pt, arc=2pt,
  left=7pt, right=7pt, top=5pt, bottom=5pt,
  fontupper=\small,
  coltitle=decisiongreen!75!black, fonttitle=\bfseries\sffamily\footnotesize,
  title={DERIVATION}
}

\newtcolorbox{boundarybox}{
  enhanced, breakable,
  colback=boundaryred!4, colframe=boundaryred!55, boxrule=0.4pt, arc=2pt,
  left=7pt, right=7pt, top=4pt, bottom=4pt,
  fontupper=\small,
  coltitle=boundaryred!75!black, fonttitle=\bfseries\sffamily\footnotesize,
  title={SCOPE BOUNDARY}
}

\title{\bfseries Before the PAIN Equation:\\
Deriving Security-Requirements Ceilings\\
from Intended Federal Information Types}
\author{Matthew Venne \\ \small Chief Technology Officer, stackArmor}
\date{July 2026\\\small Version 1.0}

\begin{document}
\maketitle

\begin{abstract}
\noindent
Potential Agency Impact (PAIN) calculations require a Confidentiality, Integrity,
and Availability Security Requirements vector for the affected asset. An asset
archetype can provide a consistent architectural vector, but it cannot by itself
answer a different question: which federal consequences are possible in a
particular agency's intended use of a Cloud Service Offering (CSO)?

This paper specifies the missing upstream derivation. It begins with the federal
information types the CSO is designed and authorized to process and the information
types an agency intends and is authorized to place in that offering. Confirmed
overlap is mapped to the provisional objective values and special factors in NIST
SP~800-60. The applicable values are aggregated independently for Confidentiality,
Integrity, and Availability, preserving the full dimensional vector required by
FIPS~199. That vector becomes a Security Requirements ceiling. The effective
requirement for an asset is the per-objective minimum of its system-relative
archetype requirement and the intended-use ceiling.

The separation matters. Archetypes remain reusable across agency deployments because
they describe the role of a component inside the CSO, not every possible customer
data type. The ceiling supplies the agency-specific consequence context without
turning a system's overall FIPS~199 impact label into an all-equal vector. Generic
free text, file upload, or integration capability does not authorize hypothetical
information types. Multi-agency assets aggregate only the definite agencies they can
affect. The assessment measures adverse effect on federal operations, federal
assets, and individuals; a CSP's separate brand or corporate risk does not silently
raise agency impact.

This paper ends at the effective Security Requirements vector. It does not choose a
PAIN aggregation formula, word threshold, or N-rating. Those downstream policy
choices belong to the PAIN method and its calibration.
\end{abstract}

\tableofcontents

\section{Status, scope, and relationship to PAIN}

This document is informational. It does not define a FedRAMP requirement, replace an
agency categorization, or make NIST SP~800-60 recommendations mandatory. It proposes
a deterministic bridge from federal information categorization to the
Security Requirements inputs consumed by \painpaper{}.

The key words \rfckw{must}, \rfckw{must not}, \rfckw{should},
\rfckw{should not}, and \rfckw{may} describe this proposed method as in RFC~2119.
They do not create obligations under FedRAMP, FIPS~199, or NIST guidance.

\begin{description}[leftmargin=1.6em,style=nextline]
\item[Cloud Service Offering (CSO)] The bounded cloud service whose intended federal
  use is being assessed.
\item[Information type] A category of federal information identified under NIST
  SP~800-60, with provisional Confidentiality, Integrity, and Availability impact
  recommendations.
\item[Applied profile] The objective vector retained after an agency confirms that
  an information type applies and evaluates its special factors and actual use.
\item[Archetype requirements] The uncapped CR/IR/AR vector assigned from an asset's
  role and consequence inside the CSO architecture.
\item[Security-requirements ceiling] The greatest C/I/A consequence supported by
  the confirmed intended information use in the asset's agency scope.
\item[Effective requirements] The per-objective minimum of the archetype
  requirements and the ceiling. This is the CR/IR/AR vector handed to PAIN.
\end{description}

\begin{boundarybox}
The paper stops at effective CR/IR/AR. Whether a particular effective vector can
produce Minimal, Narrow, Disruptive, or Debilitating under a selected PAIN formula
is a property of that downstream formula and its governed thresholds. It is not a
property established by NIST SP~800-60 or FIPS~199 alone.
\end{boundarybox}

\section{The missing first step}

FedRAMP requires providers to evaluate vulnerabilities in the context of the CSO
and estimate potential impact on government customers. CVSS Environmental Security
Requirements supply a standard vocabulary for stating how strongly the affected
environment requires confidentiality, integrity, and availability. Neither source,
however, automatically provides the three values for a particular asset in a
particular agency deployment.

The existing PAIN method recommends asset archetypes: governed role descriptions
such as system-of-record database, identity service, orchestrator, message broker,
or disposable cache. This is preferable to asking each asset owner to declare an
unexplained value. Yet one vector is being asked to carry two different kinds of
information:

\begin{enumerate}[leftmargin=1.7em]
\item \textbf{Architectural consequence:} what role does this component perform
  inside the CSO, and how much can its compromise affect the offering?
\item \textbf{Federal mission consequence:} what information and government use
  are actually in scope for the affected agency or agencies?
\end{enumerate}

Combining those questions inside every archetype makes the catalog customer-specific.
A database archetype may become High for one agency, Moderate for another, and then
change again when the CSO adds a customer. The resulting catalog is difficult to
reuse and difficult to audit.

This method separates the questions. The archetype holds the first answer. A
Security Requirements ceiling derived from intended federal information use holds
the second. PAIN receives their per-objective intersection.

\begin{plain}
A primary database can be architecturally important even when the current federal
use is only Moderate. Keep the database role stable. Derive the agency's
three-dimensional ceiling separately. The database keeps its identity, while its
effective PAIN requirements reflect what is actually at stake for that agency.
\end{plain}

\section{Risk has a point of view}

FIPS~199 defines impact through adverse effects on organizational operations,
organizational assets, and individuals. NIST SP~800-60 applies that federal frame
to information types and agency missions. FedRAMP's PAIN vocabulary likewise asks
about effects on agencies using the cloud service.

Risk is therefore not an intrinsic property of a byte sequence. The same
information can carry different consequences for different stakeholders:

\begin{table}[H]
\centering
\small
\renewcommand{\arraystretch}{1.2}
\begin{tabularx}{\textwidth}{@{}P{0.21\textwidth}YY@{}}
\toprule
\textbf{Stakeholder} & \textbf{Question} & \textbf{Treatment in this method} \\
\midrule
Federal agency &
Would disclosure, corruption, or loss cause limited, serious, or severe or
catastrophic adverse effect on the mission, federal assets, or individuals? &
This is the impact measured by the ceiling and handed to PAIN. \\
CSP &
Would the same event harm brand, market position, revenue, contractual posture, or
corporate reputation? &
Important enterprise risk, but not automatically federal agency impact. Include it
only where it creates a documented federal consequence. \\
Other customer &
Would the event harm a non-federal customer's operations or obligations? &
Outside federal PAIN unless that consequence also affects an agency in the assessed
scope. \\
\bottomrule
\end{tabularx}
\caption{Stakeholder consequences are related but not interchangeable.}
\end{table}

A Low federal Confidentiality objective does not mean disclosure is desirable or
costless. It means the agency has determined that unauthorized disclosure would
have only limited adverse effect under the FIPS~199 frame. A CSP may rationally
assign greater internal importance to avoiding the same disclosure. That separate
judgment belongs in the CSP's enterprise risk process rather than being smuggled
into PAIN as if the government made it.

\section{NIST SP 800-60 as provisional evidence}

NIST SP~800-60 Volume~II Revision~1 contains a catalog of management-and-support
and mission-based information types. Each scored record supplies a provisional
C/I/A profile, objective-level rationale, special factors, and a recommendation.
The values are starting points. The agency identifies applicable information
types, evaluates actual use, applies special factors, considers aggregation and
system context, and documents adjustments.

\subsection{What the catalog does and does not decide}

For an information type $t$, write its provisional profile as
\[
p_t=(p_t^C,p_t^I,p_t^A), \qquad p_t^o\in\{L,M,H\}.
\]
After the agency evaluates actual use and special factors, the applied profile is
\[
q_{a,t}=(q_{a,t}^C,q_{a,t}^I,q_{a,t}^A).
\]

The method \rfckw{must} preserve both values and explain every difference. A
provisional High is not automatic proof that the agency's use is High. A
provisional Low does not prevent a documented special factor or aggregate from
raising the applied objective. A candidate information type does not participate
in the math until its applicability is confirmed.

The catalog does not cover classified information or national-security systems.
Those uses remain outside this method.

\subsection{High-impact paths in the catalog}

A source-pinned conversion of Volume~II contains 170 records: 168 scored
information types and two delivery mechanisms without standalone profiles. Eleven
of the 168 scored types begin with at least one provisional High objective:

\begin{table}[H]
\centering
\scriptsize
\renewcommand{\arraystretch}{1.12}
\begin{tabularx}{\textwidth}{@{}P{0.10\textwidth}YP{0.19\textwidth}@{}}
\toprule
\textbf{ID} & \textbf{Information type} & \textbf{Provisional vector} \\
\midrule
D.2.2 & Key Asset and Critical Infrastructure Protection & C:H/I:H/A:H \\
D.2.3 & Catastrophic Defense & C:H/I:H/A:H \\
D.2.4 & EOP Executive Functions & C:H/I:M/A:H \\
D.3 & Intelligence Operations (Unclassified Domestic Intelligence) & C:H/I:H/A:H \\
D.4.1 & Disaster Monitoring and Prediction & C:L/I:H/A:H \\
D.4.4 & Emergency Response & C:L/I:H/A:H \\
D.5.1 & Foreign Affairs & C:H/I:H/A:M \\
D.5.3 & Global Trade & C:H/I:H/A:H \\
D.11.4 & Space Operations & C:L/I:H/A:H \\
D.14.4 & Health Care Delivery Services & C:L/I:H/A:L \\
D.17.2 & Legal Defense & C:M/I:H/A:L \\
\bottomrule
\end{tabularx}
\caption{The eleven information types with a provisional High objective.}
\end{table}

Only four provisionally begin at H/H/H. Another 94 scored records contain
objective-specific special-factor language that explicitly permits High under
stated conditions. Sixty-three have no explicit High path in either their
provisional profile or special factors.

These counts require careful interpretation. Provisional High is uncommon.
Conditional High is available across a broader part of the catalog, but a
conditional path does not make the information type High by itself. Common
conditions concern life safety, primary mission failure, critical infrastructure,
time-critical emergency or operational decisions, catastrophic market disruption,
and disclosure or corruption that enables catastrophic harm. Each requires an
affirmative fact about the agency's intended use.

\begin{boundarybox}
The catalog supports ``provisional High is rare'' and ``conditional High requires
specific facts.'' It does not support the claim that only a small fraction of
information types have any conceivable High path: 105 of 168 have either a
provisional or explicit conditional path. Nor does the catalog establish how many
agencies actually meet those conditions.
\end{boundarybox}

\section{Preserving the FIPS 199 dimensions}

\subsection{Information-type profiles}

FIPS~199 expresses the security category of information as a tuple:
\[
\mathrm{SC}_{t}
=\{(C,q_t^C),(I,q_t^I),(A,q_t^A)\}.
\]
The dimensions are not interchangeable. An information type may be High for
Integrity and Low for Confidentiality and Availability. That asymmetry is evidence,
not noise.

\subsection{System aggregation}

For the set of applicable information types $T$ in a system, FIPS~199 takes the
high-water mark independently for each objective:
\begin{align}
\mathrm{SC}_{\mathrm{system}}(C)&=\max_{t\in T} q_t^C,\\
\mathrm{SC}_{\mathrm{system}}(I)&=\max_{t\in T} q_t^I,\\
\mathrm{SC}_{\mathrm{system}}(A)&=\max_{t\in T} q_t^A.
\end{align}
NIST SP~800-60 then calls for review of aggregation, critical system functionality,
and other system factors, with documented adjustment where warranted.

The system's overall impact level is the highest of the three resulting objectives:
\[
\mathrm{Impact}_{\mathrm{system}}
=\max\{\mathrm{SC}_{\mathrm{system}}(C),
       \mathrm{SC}_{\mathrm{system}}(I),
       \mathrm{SC}_{\mathrm{system}}(A)\}.
\]

Therefore any applied High objective makes the information system High overall. A
genuinely Moderate system cannot contain a High system objective. But the overall
word \emph{High} does not replace the vector. If the category is H/M/L, PAIN needs
H/M/L, not H/H/H.

\begin{modelbox}
\textbf{Use the high-water mark twice, for two different outputs.}

\begin{enumerate}[leftmargin=1.5em,nosep]
\item Take a maximum per objective to retain the system vector.
\item Take the maximum of that vector to name the overall FIPS~199 impact level.
\end{enumerate}

The overall label is useful for authorization and reporting. The dimensional vector
is the input needed for an objective-sensitive PAIN calculation.
\end{modelbox}

\subsection{Certification Class is not the vector}

FedRAMP's 2026 terminology uses Certification Classes A--D for package
specifications. Current FedRAMP guidance describes those classes as levels of
assurance that can support agency use cases; it does not make the Class letter a
substitute for an agency's FIPS~199 C/I/A category. In particular, the guidance
states that some Class~C packages may support some High security objectives.

This method therefore records Certification Class separately. Class can corroborate
an intended-use assessment and selects FedRAMP processes outside this paper, but it
does not flatten the ceiling. A conflict between the agency's category, intended use,
and the CSO's authorized boundary is surfaced for review rather than resolved by
silently changing the data.

\section{Intended use, not hypothetical capability}

\subsection{Two sets}

Let:

\begin{description}[leftmargin=2.2em,style=nextline]
\item[$T_{\mathrm{CSO}}$] Information types the offering is designed and authorized
  to process, including applicable service-generated records needed to operate the
  offering.
\item[$T_a$] Information types agency $a$ intends and is authorized to process
  through the offering.
\end{description}

The applicable set for that agency and offering is:
\begin{equation}
T_{a,\mathrm{CSO}}=T_{\mathrm{CSO}}\cap T_a.
\label{eq:type-intersection}
\end{equation}

Each entry is more than an identifier. It carries the agency-confirmed applied
profile and the special-factor rationale for this use.

\subsection{Generic input mechanisms}

A text box can technically accept a medical narrative, criminal-justice record,
tax return, trade secret, grocery list, or poem. That technical fact does not place
all corresponding information types inside the intended boundary. The same is true
of file uploads, object storage, email ingestion, and general APIs.

The CSO documents what it is designed and authorized to do. The agency documents
what it intends and is authorized to put there. Only the confirmed overlap enters
Equation~\eqref{eq:type-intersection}. Prohibited information remains prohibited.
Accidental or unauthorized data spillage is managed as spillage, not normalized
into the ordinary categorization merely because the user interface could carry it.

This rule prevents generic product capability from becoming an automatic High
designation. It also prevents the opposite error: a documented free-text,
attachment, or integration workflow that is intended to carry a sensitive
information type cannot be ignored just because the product has a generic label.

\subsection{System-generated and control information}

The intersection is not limited to agency-uploaded content. A CSO can generate
audit records, identity material, security telemetry, configurations, decisions,
transactions, or derived aggregates that require categorization. NIST SP~800-60
also directs reviewers to consider critical system functionality and security
policy enforcement information.

The assessment therefore records:

\begin{itemize}[leftmargin=1.5em]
\item agency-supplied intended information;
\item service-generated information;
\item derived or aggregated information;
\item trusted decisions and actions performed by the CSO;
\item documented system-level adjustments.
\end{itemize}

These facts can alter an applied profile. They do not erase the requirement to
identify the objective and rationale being changed.

\section{From applied information types to a ceiling}

\subsection{One agency}

For agency $a$, objective $o\in\{C,I,A\}$, and confirmed applied profile
$q_{a,t}$, define:
\begin{equation}
E_a(o)=\max_{t\in T_{a,\mathrm{CSO}}}q_{a,t}^o.
\label{eq:agency-ceiling}
\end{equation}

Apply any documented system-level adjustment after the information-type maximum.
The result
\[
E_a=(E_a(C),E_a(I),E_a(A))
\]
is the intended-use ceiling for that agency's deployment.

If no intended information types have been confirmed, the method does not invent a
low ceiling. It records an unresolved assessment and uses a conservative,
explicitly provisional upper bound until the intended use is confirmed.

\subsection{Several agencies and affected scope}

Let $A_x$ be the definite agencies whose information or mission can be affected
through asset $x$. The asset-scoped ceiling is:
\begin{equation}
E_x(o)=\max_{a\in A_x}E_a(o).
\label{eq:asset-ceiling}
\end{equation}

This is deliberately asset-scoped:

\begin{itemize}[leftmargin=1.5em]
\item A single-tenant component receives the ceiling for its one definite agency.
\item A shared control plane receives the per-objective maximum for every definite
  agency it can affect.
\item Adding an agency to a shared component can raise, but never lower, its ceiling.
\item Adding an agency elsewhere does not change a component that cannot affect
  that agency.
\end{itemize}

Agency identity alone does not determine an objective. An agency may operate many
systems with different information types. Only its intended use of the assessed CSO
participates.

\subsection{The vector approximation}

Some implementations already maintain only two aggregate vectors:

\begin{itemize}[leftmargin=1.5em]
\item an offering or System Security Objectives vector (SSO);
\item an Agency Security Objectives vector (ASO).
\end{itemize}

They can approximate Equation~\eqref{eq:agency-ceiling} with:
\begin{equation}
\widehat{E}_a(o)=\min(\mathrm{SSO}(o),\mathrm{ASO}_a(o)).
\label{eq:vector-approx}
\end{equation}

For several agencies, take the per-objective maximum of their ASO vectors before
the minimum. This approximation is simple and safe when the aggregate vectors
faithfully describe the same intended-use set. It can be imprecise when the CSO and
agency vectors become High for unrelated information types. Two aggregate High
values prove dimensional compatibility, not type-level overlap.

The preferred audit record therefore retains confirmed information-type membership.
Equation~\eqref{eq:vector-approx} remains a documented fallback, not a claim that
the set intersection was evaluated.

\begin{plain}
The vector shortcut asks, ``How high can each side go?'' The type-aware method also
asks, ``Are they talking about the same intended use?'' Keep the shortcut where
detail is unavailable, but label it honestly.
\end{plain}

\section{Archetype requirements remain separate}

\subsection{Role, not software kind}

An archetype describes how an asset is used inside the CSO. The same software can
serve different roles:

\begin{itemize}[leftmargin=1.5em]
\item an in-memory store used as a disposable cache;
\item the same store holding durable sessions or authentication tokens;
\item the same store brokering mission-driving work;
\item the same engine acting as a system-of-record database.
\end{itemize}

Those roles justify different uncapped requirements even before the agency context
is applied. The archetype catalog is therefore system-relative: a provider derives
it from the CSO architecture, trust relationships, control functions, and
consequences of component compromise.

\subsection{Effective requirements}

Let the archetype requirements for asset $x$ be
\[
R_x=(R_x(C),R_x(I),R_x(A)).
\]
The effective requirements handed to PAIN are:
\begin{equation}
R_x^{\mathrm{eff}}(o)=\min(R_x(o),E_x(o)),
\qquad o\in\{C,I,A\}.
\label{eq:effective}
\end{equation}

The ceiling never raises an archetype. A High ceiling means only that High remains
available where the asset's architectural role supports it. Conversely, a High
archetype does not manufacture a High federal consequence when the intended-use
ceiling is Moderate.

\begin{table}[H]
\centering
\small
\renewcommand{\arraystretch}{1.15}
\begin{tabularx}{\textwidth}{@{}P{0.25\textwidth}P{0.22\textwidth}P{0.22\textwidth}Y@{}}
\toprule
\textbf{Input} & \textbf{Example} & \textbf{Scope} & \textbf{Meaning} \\
\midrule
Archetype requirements & H/H/M & Asset in the CSO & Uncapped architectural
consequence of the role \\
Intended-use ceiling & M/M/L & Agency or affected agency set & Maximum federal
consequence supported by confirmed intended use \\
Effective requirements & M/M/L & Asset in that agency scope & Per-objective
minimum supplied to PAIN \\
\bottomrule
\end{tabularx}
\caption{The three vectors remain visible in the audit chain.}
\end{table}

\subsection{Why this improves reuse}

The same CSO can retain one governed archetype catalog across several agency
deployments. The customer-specific information appears in the ceiling. A role
change still changes the archetype. An intended-use change still changes the
ceiling. Reviewers can challenge either judgment without rebuilding both.

Cross-system trust anchors, software-delivery paths onto agency devices, and shared
CSP infrastructure require particular care because their consequences may extend
beyond the information stored in the assessed component. The correct response is
to model their architectural role and affected scope explicitly, not to make a
data-only archetype and then hide an exception.

\section{Worked derivations}

\subsection{Moderate intended use on a high-role database}

A CSO's system-of-record database has archetype requirements H/H/M. One agency's
confirmed intended information types aggregate to M/M/L. Therefore:
\[
R_{\mathrm{db}}^{\mathrm{eff}}
=\min((H,H,M),(M,M,L))=(M,M,L).
\]

The database remains the system-of-record archetype. No customer-specific variant
is created. If a later agency uses the same CSO differently, its ceiling changes
without rewriting the database role.

\subsection{High confidentiality without H/H/H flattening}

An agency's applied information types and special factors aggregate to H/M/L. The
overall FIPS~199 impact level is High because one objective is High. The ceiling
remains H/M/L.

For a secrets service with H/H/H archetype requirements:
\[
R_{\mathrm{secrets}}^{\mathrm{eff}}
=\min((H,H,H),(H,M,L))=(H,M,L).
\]

Turning the overall High label into H/H/H would discard the agency's dimensional
judgment and overstate Integrity and Availability.

\subsection{A shared asset across two agencies}

Agency~1 has ceiling M/M/L. Agency~2 has ceiling L/H/M. A component dedicated to
Agency~1 retains M/M/L. A component dedicated to Agency~2 retains L/H/M. A shared
control plane that can affect both receives:
\[
E_{\mathrm{shared}}
=\max((M,M,L),(L,H,M))=(M,H,M).
\]

Multi-agency scope increases only the objectives supported by a definite affected
agency. It does not automatically force H/H/H.

\subsection{A generic text field}

A project-management CSO provides comments and attachments. Its documented
intended use covers project plans, task records, and approved supporting artifacts.
The agency has prohibited health, tax, criminal-justice, and classified content.

The text field does not add every NIST information type to
$T_{\mathrm{CSO}}$. Confirmed project and planning types enter the ceiling.
Prohibited content remains governed by data-use policy and spillage response.
If the agency later authorizes a new intended record type, the type inventory and
ceiling are reviewed as a change.

\subsection{Government impact and CSP reputation}

An information type has applied C:L because disclosure would have limited adverse
effect on the agency's operations. The CSP believes any public incident involving
the record would materially harm its brand. The ceiling retains C:L unless the CSP
can show how that reputational event creates a greater adverse effect on the agency,
federal assets, or affected individuals. The CSP can and should manage its separate
corporate exposure through enterprise risk controls.

\section{Cross-CSP composition and the hyperscaler boundary}

\subsection{The knowledge problem}

An upstream CSP may support many downstream CSOs without knowing which agencies
ultimately depend on each service or which information types those agencies process.
Requiring a hyperscaler to collect every downstream agency customer list is
unlikely to be practical and may conflict with contractual boundaries.

The missing information does not prove a low ceiling. It creates uncertainty about
the dimensional aggregate.

\subsection{Privacy-preserving vector attestations}

A downstream CSO could share a bounded assertion without naming its agency
customers or detailed information types:

\begin{itemize}[leftmargin=1.5em]
\item downstream CSO identifier;
\item authorization or categorization boundary;
\item maximum required C/I/A ceiling for the dependency;
\item whether the dependency is single-customer or shared;
\item evidence status, effective date, and expiry;
\item attesting party and signature.
\end{itemize}

The upstream CSP can take the per-objective maximum of current assertions. This is
less precise than type-level evidence but more informative than customer identities
alone.

\subsection{A non-normative saturation approximation}

Where no downstream vector assertions exist and the upstream service supports a
broad, heterogeneous population of confirmed downstream CSOs, the dimensional
maximum will tend toward the upper envelope allowed by the upstream boundary. A
reasonable conservative approximation is:

\[
E_{\mathrm{downstream}}\approx
\begin{cases}
(M,M,M), & \text{Moderate-impact upstream system},\\
(H,H,H), & \text{High-impact upstream system}.
\end{cases}
\]

This is an uncertainty policy, not a NIST or FedRAMP derivation. The trigger cannot
be ``many'' in undocumented prose. An adopter must govern the dependency-count or
diversity threshold, evidence window, confidence, and review cycle. The
approximation never authorizes a use outside the upstream boundary: if a downstream
use actually requires High while the upstream information system is genuinely
Moderate, the result is a boundary mismatch, not permission to process the use.

\begin{boundarybox}
This paper identifies the cross-CSP composition problem but does not standardize a
saturation threshold or attestation protocol. The approximation is suitable as a
disclosed conservative fallback. A complete supply-chain profile is future work.
\end{boundarybox}

\section{Governance and audit record}

A deterministic equation does not make uncertain inputs certain. A conforming
assessment preserves evidence and ownership for each judgment.

\subsection{Minimum record}

The record \rfckw{should} contain:

\begin{itemize}[leftmargin=1.5em]
\item CSO identity, boundary, designed functions, and service-generated data;
\item confirmed intended and prohibited information types;
\item per-agency information types, relationship (definite or target), and
  governing evidence;
\item each provisional and applied C/I/A profile;
\item special factors considered and objective-specific adjustment rationale;
\item system-level aggregation or critical-function adjustments;
\item the per-agency and asset-scoped ceilings;
\item the overall FIPS~199 impact label, stored separately from the vector;
\item archetype requirements, effective requirements, and every capped objective;
\item multi-agency affected-scope basis;
\item confidence, assumptions, manual-review items, effective date, and expiry;
\item any vector approximation or downstream-CSP saturation approximation used.
\end{itemize}

Operator attestations remain separate from analyst or automated inferences.
Confidence measures evidence quality and never lowers an objective. When two
adjacent values remain credible, select the higher value provisionally and state
what evidence would justify lowering it.

\subsection{Change triggers}

Reassess the ceiling when:

\begin{itemize}[leftmargin=1.5em]
\item an agency adds or removes an intended information type;
\item a special factor becomes applicable or ceases to apply;
\item a CSO adds a new ingestion path, generated record, trusted decision, or
  system function;
\item tenancy or shared-control-plane scope changes;
\item aggregation creates a new system-level consequence;
\item the governing categorization or authorization boundary changes;
\item a downstream-CSP assertion expires or the approximation threshold is crossed.
\end{itemize}

Archetype changes follow different triggers: component role, trust, privilege,
dependency, or architectural consequence. Keeping the triggers separate is one of
the model's principal governance benefits.

\section{Handoff contract}

The output of this paper is not a PAIN word or N-rating. It is an auditable,
three-dimensional input:

\begin{modelbox}
For each affected asset $x$:
\[
\boxed{
R_x^{\mathrm{eff}}(o)
=\min\left(
R_x(o),
\max_{a\in A_x}
\max_{t\in T_{\mathrm{CSO}}\cap T_a}
q_{a,t}^o
\right)
}
\]
for $o\in\{C,I,A\}$, with documented system-level adjustments applied where
required.
\end{modelbox}

A transport may encode the resulting vector as:

\begin{verbatim}
CR:M/IR:H/AR:L
\end{verbatim}

or:

\begin{verbatim}
cr-m_ir-h_ar-l
\end{verbatim}

The handoff must also identify the uncapped archetype vector, ceiling vector,
ceiling source and scope, and whether any objective was capped. The downstream PAIN
method then combines this effective requirements vector with vulnerability impact.

\section{Limitations and prohibited overclaims}

\begin{itemize}[leftmargin=1.5em]
\item NIST SP~800-60 profiles are provisional recommendations. They do not replace
  agency categorization, governing sources, or actual-use evidence.
\item Classified information and national-security systems are outside the
  publication's scope.
\item Information-type intersection does not eliminate system-level aggregation,
  critical-function, control-information, or extenuating-circumstance review.
\item The vector approximation can conceal mismatched information types and must be
  labeled as an approximation.
\item The hyperscaler saturation rule is conservative policy, not a measured
  probability or a federal standard.
\item Certification Class does not replace the FIPS~199 category or the intended-use
  vector.
\item The ceiling is a cap, not a floor. It cannot make a low-role asset High.
\item This paper does not prove that any effective vector produces a particular
  PAIN word. Such a result depends on the downstream PAIN formula and thresholds.
\item The catalog counts describe text and provisional profiles, not the prevalence
  of High-impact agency deployments.
\end{itemize}

\section{Conclusion}

The first input to an agency-specific PAIN calculation should be the federal
consequence that is actually possible in the assessed use. NIST SP~800-60 supplies
provisional, objective-specific evidence; FIPS~199 supplies dimensional aggregation
and the overall system impact rule. Neither calls for replacing the resulting
vector with one repeated High, Moderate, or Low label.

Separating the intended-use ceiling from the asset archetype makes both judgments
clearer. The archetype records what the component does in the CSO. The ceiling
records what the definite agency use puts at stake. Their per-objective minimum
becomes the effective CR/IR/AR vector handed to PAIN. Multi-agency components
aggregate only the agencies they can affect. Generic technical capability does not
authorize hypothetical data. CSP corporate consequences remain important but do
not silently redefine government impact.

The result is a reusable architectural classification joined to an agency-specific,
source-traceable federal consequence vector. The PAIN equation can then begin with
its most consequential assumption made explicit.

\clearpage
\section*{References}
\addcontentsline{toc}{section}{References}

{\small
\begin{itemize}[leftmargin=3.2em,itemsep=0.35em,topsep=0.4em]
\item[{[FIPS199]}] National Institute of Standards and Technology,
  \emph{Standards for Security Categorization of Federal Information and
  Information Systems}, FIPS PUB 199, February 2004.
  \url{https://doi.org/10.6028/NIST.FIPS.199}.
\item[{[SP800-60-1]}] K.~Stine, R.~Kissel, W.~Barker, J.~Fahlsing, and J.~Gulick,
  \emph{Guide for Mapping Types of Information and Information Systems to Security
  Categories, Volume I}, NIST SP 800-60 Volume I Revision 1, August 2008.
  \url{https://doi.org/10.6028/NIST.SP.800-60v1r1}.
\item[{[SP800-60-2]}] K.~Stine, R.~Kissel, W.~Barker, A.~Lee, J.~Fahlsing, and
  J.~Gulick, \emph{Guide for Mapping Types of Information and Information Systems
  to Security Categories, Volume II}, NIST SP 800-60 Volume II Revision 1,
  August 2008.
  \url{https://doi.org/10.6028/NIST.SP.800-60v2r1}.
\item[{[CVSS31]}] FIRST, \emph{Common Vulnerability Scoring System v3.1:
  Specification Document}, Environmental Security Requirements.
  \url{https://www.first.org/cvss/v3.1/specification-document}.
\item[{[FEDRAMP-VER]}] FedRAMP, \emph{Vulnerability Evaluation and Reporting},
  Consolidated Rules for 2026, rule VER-EVA-EPA, launched June 24, 2026.
  \url{https://www.fedramp.gov/2026/reference/rev5/b/vulnerability-evaluation-and-reporting/}.
\item[{[FEDRAMP-PAIN]}] FedRAMP, \emph{Initial Outcome of RFC-0031: Updated
  Incident Communications Procedures}, Notice 0012, PAIN and customer-effect
  terminology.
  \url{https://www.fedramp.gov/notices/0012/}.
\item[{[FEDRAMP-CLASS]}] FedRAMP, \emph{Choosing a Certification Class},
  Consolidated Rules for 2026.
  \url{https://www.fedramp.gov/preview/2026/providers/start/class/}.
\item[{[PAIN]}] M.~Venne, \emph{A Deterministic, CVSS--Environmental Method for
  FedRAMP VDR/VER Vulnerability Prioritization}, stackArmor, July 2026.
  \url{https://stackarmor.github.io/rfc-fedramp-vdr/vdr-pain-cvss.pdf}.
\end{itemize}
}

\subsection*{Reproducibility record}

The catalog analysis uses the source-pinned NIST SP~800-60 Volume~II Revision~1
conversion distributed with \code{trivy-plugin-vdr-skills}. A conforming run must
return:

\begin{center}
\small
\begin{tabular}{@{}lr@{}}
\toprule
\textbf{Check} & \textbf{Expected result} \\
\midrule
Total catalog records & 170 \\
Scored information types & 168 \\
Non-scoring delivery mechanisms & 2 \\
Provisional profile contains at least one High & 11 \\
Provisional H/H/H & 4 \\
Additional records with explicit conditional High language & 94 \\
Unique provisional or conditional High path & 105 \\
No explicit provisional or conditional High path & 63 \\
\bottomrule
\end{tabular}
\end{center}

The conditional count includes a record only when objective-specific
\code{specialFactors} text expressly permits a High impact level. A conditional
record contributes no High objective to an assessed system unless the agency
confirms that the stated factor applies.

\end{document}
